\chapter*{Fondements, objectifs et résultats structurels}
\addcontentsline{toc}{chapter}{Fondements, objectifs et résultats structurels}
\markboth{Fondements et résultats structurels}{Fondements et résultats structurels}
\label{hmt:sistematica:presentacion}

L’Holographie Modulaire Triadique étudie la production de structures numériques, géométriques et opératoires à partir d’une arithmétique discrète orientée. Son objet initial est une grille construite avec les entiers de un à neuf. L’addition et la multiplication déterminent deux évaluations de chaque position ; les réductions modulaires conservent les restes et les quotients ; une unité ternaire sélectionne les régimes opératoires, et un système de transport ordonne leurs transitions et leur mémoire. L’itération compatible de ces opérations constitue l’état enrichi et son prolongement coinductif.

La thèse directrice attribue au nombre une structure antérieure à son évaluation scalaire. La position, l’orientation, la retenue, les incidences et la chronologie des opérations demeurent dans l’état produisant la valeur. La réalisation archimédienne s’obtient par contraction de cylindres compatibles ; la réalisation conjuguée conserve l’incidence de frontière. Toutes deux sont des lectures d’une même construction. En ce sens précis, la droite réelle apparaît comme résultat de la double projection.

Ce travail développe systématiquement cette généalogie. Ses résultats principaux s’organisent en quatre niveaux : la génération régionale coinductive des constantes ; l’articulation entre évaluation et incidence exceptionnelle ; les lois de conservation et de récupération de l’information ; et les réalisations de Mécanique Dimensionnelle. La clôture cardinale du continu engendré et les réalisations terminales reposent sur cette architecture, avec les domaines, opérateurs et hypothèses établis dans leurs théorèmes.

\section*{Définition opératoire du noyau formel}

\emph{All Possible Paths}, abrégé APP, désigne le support arithmétique des trajectoires sur neuf lignes et neuf colonnes. Sa matrice principale se construit en multipliant les indices de ligne et de colonne puis en appliquant la racine numérique positive. Si \(D_9=\{1,\ldots,9\}\), ses positions sont \((i,j)\in D_9^2\), et l’entrée multiplicative est
\[
 \Pi(i,j)=1+((ij-1)\bmod9).
\]
Le feuillet additif évalue \(i+j\) sur les mêmes positions. Sur chaque feuillet, le couple formé par le reste positif et le quotient restitue l’entier : \(n=\rho_9(n)+9q_9(n)\). Le symbole \(9\) représente la classe nulle dans la section positive. Cette convention rend visible la frontière modulaire et conserve son relèvement arithmétique \eqref{eq:app-restitution}.

Le TRIT est l’unité ternaire orientée, de valeurs \(-1,0,+1\). Il agit comme lecteur de classes opératoires et comme sélecteur de régime. Sa décomposition arithmétique conserve le couple \(n=r_3(n)+3c_3(n)\), où \(r_3(n)\) est le représentant équilibré et \(c_3(n)\) son quotient. L’orientation et la sélection de mode acquièrent ensuite des réalisations quadratiques et temporelles, définies sur leurs supports respectifs. La composition entre les modules trois et neuf transporte les deux quotients ; l’identité \eqref{trit:eq:compatibilidad39} détermine exactement ce transport.

Le \emph{Topological Phase Kernel}, TPK, est le noyau topologique de phase. Il réunit les opérateurs de sélection, de déplacement, d’émission, de mise à jour, de relèvement et de retour. Une transition reçoit un état déterminé, décide l’opération admissible et conserve le registre nécessaire à sa continuation. L’état enrichi réunit ainsi la position d’APP, le feuillet actif, l’orientation tritique, les restes, les retenues, la phase, les tours et les incidences accumulées. Les définitions et les règles d’émission sont développées dans les sections~\ref{app:construccion-explicita}, \ref{trit:capitulo} et~\ref{tpk:capitulo}.

Les représentations géométriques se construisent sur ces données. Les translations d’APP déterminent un graphe de Cayley ; l’identification périodique des bords produit sa réalisation toroïdale. La structure de Kähler correspond à la réalisation différentielle munie d’une métrique et d’une structure complexe déclarées. Le bloc central \(\left(\begin{smallmatrix}7&2\\2&7\end{smallmatrix}\right)\) admet une normalisation déterminantale conservant une forme de Minkowski de dimension deux. Chaque réalisation conserve l’objet et la normalisation dont elle provient \ref{geo:kahler}--\ref{geo:lorentz}.

\section*{Génération régionale et prolongement coinductif}

Le premier résultat directeur consiste à passer du domaine fini des exécutions à une génération compatible à profondeur arbitraire. L’énumération initiale distingue \(104\,976\) configurations, \(468\) émissions distinctes et \(243\) mots ternaires admissibles parmi les \(729=3^6\) mots possibles. L’action orbitale correspondante produit \(43\) orbites. Chaque cardinal dénombre un objet défini : configurations, émissions, mots ou orbites.

Les relèvements prolongent les mots selon la suite
\[
 w_6\longrightarrow w_{12}\longrightarrow w_{18}
 \longrightarrow w_{24}\longrightarrow w_{30}
 \longrightarrow R_{36}\longrightarrow G_9.
\]
Les deux premiers relèvements utilisent \(L_0\) ; les suivants utilisent \(L_1\) et conservent le registre lors du changement de régime. La frontière \(R_{36}\) ouvre la continuation coinductive, et \(G_9\) réunit les neuf réalisations locales. L’holonomie \(\Gamma_9\) agit sur l’état enrichi. Sa réduction de phase est une lecture ultérieure.

Le calendrier conserve simultanément la phase nonadique et le registre dodécaphasique au moyen de
\[
 0\longrightarrow C_{12}\longrightarrow C_{108}
 \longrightarrow C_9\longrightarrow0.
\]
La retenue de cette extension détermine la différence entre un retour de phase et une mise à jour de l’état. Après neuf transitions, la phase initiale réapparaît et la mémoire des tours augmente. Le théorème~\ref{thm:cal-extension} et le développement de l’holonomie~\ref{chap:sucesor-102-08} précisent la loi de composition.

Les régions de \(\pi\), \(\varphi\) et \(e\) s’obtiennent par les émetteurs, relèvements, filtres et conditions de survie du même TPK. Leurs cylindres compatibles déterminent les développements à toute profondeur demandée. Les cinq régions de \(\pi\), l’axe de \(\varphi\), les continuations conjointes et l’involution spéculaire se distinguent par leurs opérations et leurs registres. Les caractérisations de Machin, de Perron--Fibonacci et des semi-groupes exponentiels interviennent comme identifications ou certificats ultérieurs des sorties construites. L’ordre de génération demeure fixé par les opérateurs régionaux.

\section*{Constante de couplage et structure fine}

La constante holographique de couplage arithmético-géométrique est le registre dodécaphasique engendré \(K\), avec son ordre de fenêtres, son origine de phase, son orientation et ses incidences. Sa détermination constitue un résultat central : l’horloge de conversion, le retour nonadique et le lecteur phase--charge produisent le descripteur régional ; les conditions d’incidence et d’hétérotypie sélectionnent le registre dans le domaine terminal déclaré. Le chapitre~\ref{chap:despliegue-selector-k} conserve séparément le descripteur, le répertoire de blocs, les conditions de sélection et l’inversion du registre signé.

La singularité mathématique de \(K\) s’exprime par les opérations qu’il articule. Sa lecture positionnelle intervient dans la clôture arithmétique ; ses incidences déterminent des directions et des projecteurs ; son orbite complète fournit une référence pour la récupération des opérateurs. La lecture rationnelle \(\kappa\), le repère orbital \(O_K\) et les opérateurs de mémoire reçoivent des définitions propres. Cette distinction permet de montrer la composition effective entre nombre, géométrie et restitution.

Les coordonnées régionales et le registre dodécaphasique concourent à la détermination de la constante de structure fine. La voie entière normalise les blocs avec leur retenue et conserve la frontière lors du prolongement de l’exécution. La voie analytique développe une fonction complète, avec récurrence de tous ses coefficients, et démontre sa compatibilité avec la précoordonnée engendrée. Le prolongement analytique est déterminé dans sa classe déclarée à partir du même état ; les deux voies expriment une quadrirelation coinductive commune \((\pi,\varphi,e,\alpha)\). L’ordre interne distingue les lectures régionales préalables, la sélection de \(K\), la clôture dodécaphasique et la représentation analytique corrélée \ref{x_reciprocidad:subsec:alpha-lector-completo-rev08}.

\section*{Double projection et incidence exceptionnelle}

La seconde thèse directrice est la construction conjointe du continu. Les prolongements compatibles constituent un système projectif ; leurs cylindres déterminent l’évaluation archimédienne. L’autre projection conserve les incidences, orientations et marques de frontière qui soutiennent les réalisations exceptionnelles. La contraction de l’intervalle de lecture et la persistance de cette incidence décrivent des propriétés simultanées du même état.

La construction spectrale, le bilan solénoïdal, l’incidence de jauge, la cohérence cohomologique et la droite déterminante s’obtiennent corrélativement à partir du système commun. Leurs opérateurs commutent avec les raffinements dans les domaines démontrés. Cette unité fournit le contenu des chapitres~\ref{chap:sucesor-102-05} et~\ref{chap:sucesor-102-06} et gouverne les transports ultérieurs.

La chaîne exceptionnelle développe des matrices de Paley, des incidences de Witt, des codes de Golay et des recollements réticulaires. Le voisin marqué de Leech est caractérisé par l’intégralité, la parité, l’unimodularité et l’exclusion des racines. Son prolongement par l’algèbre réticulaire, les oscillateurs, le cocycle et le secteur tordu conduit au module Moonshine. Le Monstre agit alors sur l’algèbre d’opérateurs de vertex obtenue. Ces résultats sont précisés dans les théorèmes~\ref{x_reciprocidad:exc:teorema-leech} et~\ref{x_reciprocidad:exc:moonshine}. La réduction orthogonale \(12\to11\to10\), la réduction isotrope \(26\to24\) et la dualité entre impulsion et enroulement sont coordonnées par leurs applications quotientées. La relation avec la théorie M conserve cette réalisation algébrique ainsi que les champs, actions et conditions supplémentaires de sa réalisation physique.

\section*{Conservation évolutive et récupération de l’information}

La conservation évolutive de l’unité se démontre au moyen d’applications conservant la mesure ou la norme et d’inverses restituant les données transportées. L’extrême et moyenne raison holographique se développe par la complétion \(1000=729+271\), le transfert entre échelles et la conservation de l’unité normalisée. La retenue, la mémoire de transport, l’incidence de frontière et l’archive opératoire possèdent des fonctions complémentaires dans cette articulation.

La loi bilatérale établit une identité valable pour deux transports isométriques entre espaces de Hilbert, en dimension finie ou infinie. Dans sa spécialisation nonadique, les deux canaux \(Q_-\) et \(Q_+\) satisfont
\[
 9\|Q_-v\|^2+8\|Q_+v\|^2
 =17\cdot73\,\|v\|^2.
\]
Les canaux normalisés constituent une isométrie ; son adjoint récupère l’état et a pour norme un. La loi identifie en outre les lecteurs qui demeurent invariants et la transformation des autres. La récupération à profondeur arbitraire comprend une borne uniforme pour la composante terminale et une série inverse convergente \ref{x_reciprocidad:lib04:teo-balance}, \ref{x_reciprocidad:lib04:teo-archivo}.

La réciprocité du registre étend cette conservation aux déformations de covolume unité, aux réseaux duaux, aux mineurs de Plücker et aux formes spectrales. La direction du flot provient du registre engendré ; les rapports régionaux agissent ensuite comme marques. Les démonstrations déterminent à la fois les invariants conservés et les orientations nécessaires pour inverser la lecture géométrique.

\section*{Réalisations de Mécanique Dimensionnelle}

La Mécanique Dimensionnelle développe les réalisations physiques dans l’ordre de leurs antécédents. La coordonnée de structure fine et les lectures régionales interviennent dans l’échelle d’action. Le déterminant local et son quotient d’ordre seize fixent la valuation produisant la décennie \(-34\) ; les sections aller et retour distinguent les deux lectures de Planck. L’angle et le contre-angle sont construits avec ces coordonnées et conservent leurs conversions entre radians et degrés.

L’incidence hexadique--octadique produit les canaux \(90/120\). Leur évaluation sur le couple angulaire détermine la fonctionnelle de Barbero--Immirzi ; le transport angulaire--Witt conserve la relation avec les coordonnées de mélange. La réponse constitutive du vide reconstruit ses canaux électriques et magnétiques, tandis que la section gravitationnelle relie l’action à l’unité aréale. La transduction de Boltzmann exprime l’énergie par décision relativement à la section thermique déclarée. Chaque constante conserve sa définition structurelle, ses unités et son opération d’évaluation.

La structure électronique se situe à l’unique centre fixe de l’atlas, \((5,5)\). Les autres réalisations matérielles s’organisent au moyen de parcours, d’orbites et de multisections. L’atlas de quatre-vingt-une cellules, les cinquante-six types de parcours et les treize multisections précèdent les lecteurs de masse. La composition \(\text{extension}\to\text{route}\to\text{registre}\to
\text{signature}\to\text{caractère}\to\text{tours}\) conserve la provenance de chaque observable. Les comparaisons métrologiques évaluent ensuite les représentations et normalisations correspondantes.

\section*{Clôture cardinale et régime des démonstrations}

La génération du continu reçoit une clôture cardinale explicite. Dans la clôture généalogique \(\mathfrak G\) définie par les règles de termes et de définissabilité du traité, la condensation situe chaque coupure rationnelle interne à une étape dénombrable. Le code \(j(r)=\omega\beta(r)+n(r)\) distingue leurs valeurs ; le codage des bons ordres dénombrables fournit l’injection inverse. Le théorème~\ref{ix:thm:decision-continuo-hmt} conclut
\[
 \mathfrak G_{\rm HMT}(h,m_\infty)
 \models 2^{\aleph_0}=\aleph_1,
\]
et quantifie la totalité des sous-ensembles internes de la droite engendrée. La structure et les paramètres de cette clôture proviennent de la généalogie décrite ; sa représentation par \(L[h,m_\infty]\) identifie ultérieurement le domaine du théorème.

L’exposé distingue les identités arithmétiques, les énumérations exhaustives de domaines finis, les théorèmes de compatibilité et les passages à la limite. Les formalisations vérifiées et les certificats numériques accompagnent leurs énoncés de quantificateurs et de dépendances précis. Les réalisations physiques ajoutent les cartes dimensionnelles, les couplages et les conditions d’observation qu’elles utilisent. Cette différenciation permet d’apprécier la force propre de chaque résultat et de suivre sa composition dans une seule construction.
